\chapter{Open problems, research directions and industry}
\label{ch5:ch:open}

The first four chapters followed one theorem: for every completely monotone $g$, the triangular lattice $\Lat$ of
covolume one has the least lower energy among locally finite planar configurations of density one
\cite{oai}. This chapter is about the edges of that theorem. Some are places where the theorem is silent and
something else must be said: divergent interactions, the strong Cohn--Kumar conjecture, the endpoint $s=2$. One is a
detail of the definitions that has now been settled (open versus closed disks). The rest are questions the theorem
does not ask: uniqueness, temperature, other dimensions, engineering.

Recall the status of the main theorem itself: it is AI-generated (an OpenAI release), kernel-checked in Lean and
accepted by Comparator in this project, and not peer reviewed \cite{oai}. Every statement below carries a tier. Most of this chapter is \tierC{} (argument, conjecture, proposal) or
\tierL{} (cited literature). That is the honest state of an open-problems chapter, and it should not be read as a
weakness. Table~\ref{ch5:tab:open-status} summarises where things stand.

\begin{table}[ht]
\centering\small
\caption{Status of the questions treated in this chapter. ``Lean'' means a kernel-checked statement in the
toolchain of Chapter~2, with axiom footprint inside $\{\texttt{propext},\texttt{Classical.choice},\texttt{Quot.sound}\}$.}
\label{ch5:tab:open-status}
\begin{tabular}{@{}p{4.6cm}p{4.4cm}p{4.6cm}@{}}
\toprule
Question & Mathematical status & Lean status\\
\midrule
Riesz $s>2$, Gaussians, every completely monotone $g$ & proved upstream \cite{oai} & kernel-checked (upstream + our D1, D1b, D2)\\
Raw energies for $0<s\le2$ & infinite for every competitor (\S\ref{ch5:sec:div}) & divergence of the lattice sum not formalised\\
Renormalised / jellium, $0\le s<2$ & proved in upstream's manuscript, \S8 \cite{oai}; implication from Petrache--Serfaty \cite{ps} & not formalised anywhere\\
Endpoint $s=2$ & open; outside upstream's statement & ---\\
Strong (per-potential) Cohn--Kumar & open & ---\\
Open vs closed disks & resolved in this project (D4) & see \S\ref{ch5:sec:disks}\\
Uniqueness among periodic configurations & open & ---\\
Non-uniqueness among all configurations & proved (D3) & kernel-checked\\
Positive temperature & open & ---\\
Dimension $3$ & no universal optimum known & ---\\
Dimensions $8$, $24$ & proved \cite{ckmrv} & not formalised here\\
\bottomrule
\end{tabular}
\end{table}

%=====================================================================
\section{The divergent regime \texorpdfstring{$0<s\le2$}{0<s<=2} and the logarithm}
\label{ch5:sec:div}

\subsection{Why the raw sums diverge}

Write $N(r)=\#\{a\in\Lat:0<|a|\le r\}$. The lattice has covolume one, so $N(r)$ is close to the area $\pi r^2$.

\begin{lemma}[Lattice point count]\label{ch5:lem:count}
There is a constant $C$ such that $|N(r)-\pi r^2|\le C(r+1)$ for all $r\ge0$.
\end{lemma}

\begin{humanproof}
\tierC{} Let $P$ be a fundamental parallelogram of $\Lat$ with $0\in P$, of area $1$ and diameter $\delta$. The
translates $a+P$, $a\in\Lat$, tile the plane. Every translate with $|a|\le r$ lies in the disk of radius $r+\delta$,
and every point of the disk of radius $r-\delta$ lies in a translate with $|a|\le r$. Comparing areas,
\[
\pi(r-\delta)^2\le N(r)+1\le\pi(r+\delta)^2\qquad(r\ge\delta),
\]
so $|N(r)-\pi r^2|\le 2\pi\delta r+\pi\delta^2+1$. For $r<\delta$ both $N(r)$ and $\pi r^2$ are bounded.
\end{humanproof}

\begin{proposition}[Divergence rates]\label{ch5:prop:divrate}
Let $r_0=\min_{a\ne0}|a|$. As $R\to\infty$,
\begin{gather*}
\sum_{0<|a|\le R}|a|^{-2}=2\pi\log R+O(1),\\
\sum_{0<|a|\le R}|a|^{-s}=\frac{2\pi}{2-s}R^{2-s}+O(R^{1-s}+\log R+1)\qquad(0<s<2).
\end{gather*}
\end{proposition}

\begin{humanproof}
\tierC{} Summation by parts (Stieltjes integration against the step function $N$) gives, for $s>0$,
\[
\sum_{0<|a|\le R}|a|^{-s}=\int_{r_0^-}^{R}r^{-s}\,\mathrm{d}N(r)
=N(R)R^{-s}+s\int_{r_0}^{R}N(r)\,r^{-s-1}\,\mathrm{d}r .
\]
Insert $N(r)=\pi r^2+E(r)$ with $|E(r)|\le C(r+1)$ from Lemma~\ref{ch5:lem:count}. For $s=2$:
\begin{gather*}
N(R)R^{-2}=\pi+O(R^{-1}),\qquad
2\int_{r_0}^{R}\pi r^2\,r^{-3}\,\mathrm{d}r=2\pi\log R-2\pi\log r_0,\\
2\int_{r_0}^{R}|E(r)|\,r^{-3}\,\mathrm{d}r\le 2C\int_{r_0}^{\infty}(r+1)r^{-3}\,\mathrm{d}r<\infty .
\end{gather*}
The error integral converges absolutely, so the sum is $2\pi\log R+O(1)$. For $0<s<2$ the main term is
$\pi R^{2-s}+s\pi\int_{r_0}^{R}r^{1-s}\,\mathrm{d}r=\pi R^{2-s}+\frac{s\pi}{2-s}(R^{2-s}-r_0^{2-s})
=\frac{2\pi}{2-s}R^{2-s}+O(1)$, and the error is $s\int_{r_0}^R C(r+1)r^{-s-1}\,\mathrm{d}r$, which is $O(R^{1-s})$ for
$s<1$, $O(\log R)$ for $s=1$ and $O(1)$ for $s>1$.
\end{humanproof}

The heuristic behind the proof is worth keeping: a lattice of density one looks, from far away, like a uniform area
density, and $\int_{1\le|x|\le R}|x|^{-s}\,\mathrm{d}x=2\pi\int_1^R r^{1-s}\,\mathrm{d}r$, which is $2\pi\log R$ for
$s=2$ and grows like $R^{2-s}$ for $s<2$. The sum diverges at infinity exactly when this integral does, namely when
$s\le d=2$.

\begin{numbox}[(illustration of Proposition~\ref{ch5:prop:divrate})]
\tierX{} The script below sums $|a|^{-2}$ over $0<|a|\le R$ for the covolume-one triangular lattice
$p(j,k)=b^{-1/2}(j+k/2,\,kb)$, $b=\sqrt3/2$, the parametrisation of Chapter~1. It was run as
\texttt{taskset -c 2,3 .venv/bin/python -I ch5\_checks.py}.
\begin{lstlisting}[language=Python]
def tri_s2_partial(R: float) -> float:
    b = math.sqrt(3) / 2
    sc = 1 / math.sqrt(b)
    K = int(R * 1.2 / (sc * b)) + 3
    j, k = np.meshgrid(np.arange(-K, K + 1), np.arange(-K, K + 1), indexing="ij")
    x = sc * (j + k / 2.0)
    y = sc * (k * b)
    r2 = x * x + y * y
    m = (r2 > 0) & (r2 <= R * R)
    return float(np.sum(1.0 / r2[m]))
\end{lstlisting}
Output:
\begin{lstlisting}
R=    50  S=27.101261  S-2pi logR=2.521295
R=   100  S=31.457323  S-2pi logR=2.522185
R=   200  S=35.808069  S-2pi logR=2.517760
R=   400  S=40.163973  S-2pi logR=2.518491
R=   800  S=44.519131  S-2pi logR=2.518477
\end{lstlisting}
The partial sums grow by about $2\pi\log 2\approx4.36$ per doubling of $R$, and $S(R)-2\pi\log R$ stays near $2.52$
over this range. The fluctuations come from the lattice-point error $E(r)$. We claim no more digits of the constant
and no closed form for it.
\end{numbox}

\subsection{What the Lean theorem says in this regime}

Upstream's Lean theorem (Chapter~2) states, for admissible $g$,
\begin{center}
\lean{latticeEnergy g} $\le$ \lean{energy g C} $\wedge$ \lean{latticeEnergy g} $=$ \lean{energy g A},
\end{center}
with both sides in $[0,\infty]$. Our \lean{riesz_admissible} shows that $t\mapsto t^{-s/2}$ is admissible for every
$s>0$ \tierA{}. For $0<s\le2$, Proposition~\ref{ch5:prop:divrate} shows that the lattice sum is $+\infty$
(the sum of nonnegative terms dominates every partial sum). The Lean inequality then forces
\lean{energy (riesz s) C} $=\top$ for every locally finite density-one competitor $C$. That is a true statement, but
the comparison it contains is $\infty\le\infty$, so it carries no information about which configuration is better.
This deduction is \tierC{}: our D1b proves \lean{latticeEnergy_riesz_lt_top} for $s>2$ \tierA{}, but the
complementary statement \lean{latticeEnergy (riesz s)} $=\top$ for $0<s\le2$ has not been formalised (it is
Project~M2 in \S\ref{ch5:sec:projects}).

The logarithm $g(t)=-\tfrac12\log t$ (that is, $-\log|x|$) is not even in the class: it is negative for $t>1$, and
\lean{AdmissiblePotential} requires $g\ge0$. So the theorem says nothing at all about the true two-dimensional
Coulomb interaction. Yet the cases $s=1$ (the electron gas, \S\ref{ch5:sec:industry}) and $s=0$ (the one-component
plasma, vortices) are the physically central ones. To say something there, the energy itself has to be redefined.

\subsection{Local integrability and the jellium energy}

The standard cure is to neutralise the points against a uniform background of opposite charge and the same density,
so that the divergent far-field contributions cancel. This needs the background to have a finite self-interaction,
and that is a local question:
\begin{equation}\label{ch5:eq:locint}
\int_{|x|<1}|x|^{-s}\,\mathrm{d}x=2\pi\int_0^1 r^{1-s}\,\mathrm{d}r=
\begin{cases}\dfrac{2\pi}{2-s},& s<2,\\[1ex] +\infty,& s\ge2.\end{cases}
\end{equation}
So $|x|^{-s}$ is locally integrable in $\R^2$ if and only if $s<2$, and the logarithm is locally integrable as well.
Put $g_s(x)=|x|^{-s}$ for $0<s<2$ and $g_0(x)=-\log|x|$. Upstream \cite[\S8, eq.~(finite-ordered-jellium)]{oai}
defines, for $N$ distinct points $a_1,\dots,a_N$ and the square $K_R=[-R/2,R/2]^2$ with $R=\sqrt N$, the
\emph{finite ordered-pair jellium energy}
\begin{equation}\label{ch5:eq:jellium}
J_{s,R}(a_1,\dots,a_N)=\sum_{i\ne j}g_s(a_i-a_j)-2\sum_{i=1}^N\int_{K_R}g_s(a_i-x)\,\mathrm{d}x
+\int_{K_R}\int_{K_R}g_s(x-x')\,\mathrm{d}x\,\mathrm{d}x',
\end{equation}
with $J=+\infty$ at collisions, and the density-one thermodynamic minimum
\[
e^{(s)}_{\mathrm{Jel}}=\lim_{N\to\infty}\frac1N\min_{a_1,\dots,a_N\in K_{\sqrt N}}J_{s,\sqrt N}(a_1,\dots,a_N).
\]
The three terms are point--point, point--background and background--background. Each of the last two is finite
precisely by~\eqref{ch5:eq:locint}. Note the convention: the sum is over \emph{ordered} pairs, so each pair is counted
twice. Physics papers often count each pair once; a factor $2$ between conventions must be checked before comparing
numbers. The existence of the limit is part of upstream's argument; it relies on scalar thermodynamic results of
Lewin--Lieb--Seiringer ($0<s<2$) and Lauritsen ($s=0$), both cited in \cite[\S8]{oai}.

\subsection{Upstream's renormalised corollary}

Upstream also uses the \emph{renormalised energy} $W_s$ of Petrache and Serfaty. It is defined through a field
(electric-field) formulation: for $0<s<2$ one passes to $\R^2\times\R$ with the weight $|y|^{\gamma}$, $\gamma=s-1$,
so that $g_s$ becomes the fundamental solution of $-\operatorname{div}(|y|^\gamma\nabla\,\cdot)$ with constant
$\kappa_s=4\pi$ ($\kappa_0=2\pi$ for the logarithm). The fields are taken in $L^p_{\mathrm{loc}}$ with
\[
1<p<\min\Bigl(2,\frac2s,\frac3{s+1}\Bigr)\qquad(0<s<2),
\]
the energy is truncated at a small radius $\eta$, averaged over large squares and the truncation removed; $W_s(C)$
is the infimum over compatible fields. We do not reproduce the full definition; it follows Petrache--Serfaty's
2017 field construction as cited in \cite[\S8]{oai}. On a torus $T=\R^2/L$ of area $N$ with distinct points
$a_1,\dots,a_N$, the periodic value $W_{s,L}$ has the closed form
\[
W_{s,L}(a_1,\dots,a_N)=\kappa_s^2\Bigl[\frac1N\sum_{i\ne j}\mathcal G_L(a_i-a_j,0)
+\lim_{x\to0}\Bigl(\mathcal G_L(x,0)-\frac{g_s(x)}{\kappa_s}\Bigr)\Bigr],
\]
where $\mathcal G_L$ is the mean-zero unit-source Green function of the extension problem \cite[\S8]{oai}. For a
periodic configuration $C=\bigcup_{i=1}^N(a_i+L)$ of density one write $V_s(C)=W_{s,L}(a_1,\dots,a_N)$.

\begin{corollary}[Planar renormalised and jellium energies; upstream, \S8]\label{ch5:cor:jellium}
\tierL{} Let $0\le s<2$. For every positive integer $n$ put $N=n^2$ and enumerate $\Lat/(n\Lat)$ by
$a^\triangle_1,\dots,a^\triangle_N$. Every configuration of $N$ distinct points in $\R^2/(n\Lat)$ satisfies
\[
W_{s,n\Lat}(a_1,\dots,a_N)\ge W_{s,n\Lat}(a^\triangle_1,\dots,a^\triangle_N).
\]
Moreover, $\Lat$ minimises $W_s$ among density-one planar configurations, and
$e^{(s)}_{\mathrm{Jel}}=\kappa_s^{-1}W_s(\Lat)$.
\end{corollary}

The engine of the proof is an identity that compares two periodic configurations \emph{with different period
lattices} through heat kernels.

\begin{lemma}[Heat comparison across period lattices; upstream, \S8]\label{ch5:lem:heat}
\tierL{} Let $0\le s<2$ and let $C$, $D$ be simple periodic configurations of density one. Put
\[
q_0=2\pi,\qquad q_s=\frac{2^{2-s}\pi\,\Gamma(1-s/2)}{\Gamma(s/2)}\ (0<s<2),\qquad
\Psi_t(x)=(4\pi t)^{-1}e^{-|x|^2/(4t)}.
\]
If $C=\bigcup_{i=1}^N(a_i+L)$ with $|\R^2/L|=N$, define
\[
P_{t,L}(x)=\sum_{v\in L}\Psi_t(x-v),\qquad
H_C(t)=\frac1N\sum_{i,j=1}^N P_{t,L}(a_i-a_j)-\Psi_t(0),
\]
and define $H_D$ with its own period lattice and point count. Then the following integral converges absolutely and
\[
V_s(C)-V_s(D)=\frac{\kappa_s q_s}{\Gamma(1-s/2)}\int_0^\infty\bigl(H_C(t)-H_D(t)\bigr)\,t^{-s/2}\,\mathrm{d}t .
\]
\end{lemma}

\paragraph{How the Gaussian theorem enters.}
\tierL{} The subtraction of $\Psi_t(0)$ removes only the coincident point. Counting periodic images in large disks
gives \cite[\S8]{oai}
\[
H_C(t)=\frac1{4\pi t}\,E_{g_t}(C),\qquad g_t(u)=e^{-u/(4t)},
\]
the Gaussian energy with parameter $\alpha=1/(4\pi t)$ in the normalisation $e^{-\pi\alpha u}$. The Gaussian case of
the main theorem gives $E_{g_t}(C)\ge E_{g_t}(\Lat)$, hence $H_C(t)\ge H_\Lat(t)$ for every $t>0$. The prefactor is
positive: $\kappa_sq_s/\Gamma(1-s/2)=\kappa_s2^{2-s}\pi/\Gamma(s/2)>0$, and $t^{-s/2}>0$. Lemma~\ref{ch5:lem:heat} with
$D=\Lat$ then gives $V_s(C)\ge V_s(\Lat)$ for every periodic $C$. Upstream passes to the infinite system by a
square-periodic approximation of the field infimum, which is why the lemma must allow \emph{different} period
lattices: the approximants are square-periodic while the minimiser is triangular. The identification of
$e_{\mathrm{Jel}}$ uses the scalar thermodynamic inputs cited above and a confinement argument based on a weighted
strong minimum principle (Fabes--Kenig--Serapioni, cited in \cite[\S8]{oai}), which needs $-1<\gamma<1$.

The structure is the same as in the Bernstein step of Chapter~2: a nonnegative superposition of Gaussian
inequalities. Here the weight is $t^{-s/2}\,\mathrm{d}t$, and the superposition is applied to \emph{differences}, so
the divergent parts cancel before anything is integrated. This is why the renormalised problem follows from the
Gaussian theorem even though the raw energies are infinite.

Petrache and Serfaty had earlier shown \tierL{} that universal optimality of the triangular lattice, in the form
conjectured by Cohn and Kumar, implies the crystallisation results for the Coulomb and Riesz renormalised energies
\cite[Theorem~2]{ps}. Upstream's \S8 carries out that transfer for its stated field normalisation and order of
limits, starting from its own Gaussian theorem rather than from the conjecture.

\begin{warnbox}[(what is not formalised)]
\tierC{} \emph{Nothing} in this section is in Lean. No Lean file in openai/math mentions jellium or renormalised
energy (checked by grep on the pinned tree, as recorded in the extension plan of this project). Our own Lean work
(Chapter~2) covers only summable interactions. Corollary~\ref{ch5:cor:jellium} and Lemma~\ref{ch5:lem:heat} are \tierL{}:
they are proved in an AI-generated manuscript that has not been peer reviewed \cite{oai}. The field formulation,
the thermodynamic limits and the confinement argument would each be a substantial formalisation project.
\end{warnbox}

\subsection{A Lean roadmap}

The extension plan of this project (\texttt{docs/EXTENSION\_RENORMALIZED\_JELLIUM\_2026-10-11.md}) splits off
the pieces that look reachable. Its priors are subjective, recorded in the plan before any attempt; they are not
results.

\begin{center}\small
\begin{tabular}{@{}lp{6.8cm}cp{3.6cm}@{}}
\toprule
\# & Target & Prior & Kill criterion\\
\midrule
E3 & For every periodic density-one $C$ and $t>0$: $H_C(t)\ge H_\Lat(t)$, derived from upstream's Gaussian theorem
  & $\sim30\%$ & lower energy of a periodic $C$ cannot be identified with its periodic Gaussian sum\\
E4 & $t^{-s/2}=\Gamma(s/2)^{-1}\int_0^\infty u^{s/2-1}e^{-tu}\,\mathrm{d}u$, and the positivity transfer: if
  $H_C\ge H_\Lat$ pointwise then the $t^{-s/2}$-weighted integrals compare & $\sim70\%$ & Mathlib's Gamma-integral
  API does not reach the needed form\\
E5 & Definition of $J_{s,R}$ in~\eqref{ch5:eq:jellium}, with a non-vacuity check (finite at distinct points) &
  $\sim60\%$ & the double integral needs an integrability lemma Mathlib lacks for $s$ near $2$\\
E6 & Endpoint $s=2$ & --- & outside upstream's statement; not attempted\\
\bottomrule
\end{tabular}
\end{center}

None of E3--E5 had been started when this chapter was written. If E3 and E4 both succeed, they give a kernel-checked
version of the second half of the heat-comparison argument (the inequality $V_s(C)\ge V_s(\Lat)$, \emph{given} the
identity of Lemma~\ref{ch5:lem:heat}). They would not give the field formulation, the identity itself or the
thermodynamic limit. Even the easiest step, E4, is low value on its own; its point is to be the first brick.

\subsection{A numerical check at \texorpdfstring{$s=1$}{s=1}}

\begin{numbox}[(2D Wigner crystal, $s=1$ jellium)]
\tierX{} On the branch \texttt{physics/jellium-wigner-2026-10-11} of the rusty-SUNDIALS crate
\texttt{qf-yukawa-crystal} \cite{rusty} (preregistration e706f85, code 3bc55fa), the $s=1$ jellium energy per
particle at density one was computed by Ewald summation. The triangular value is $e_{\mathrm{tri}}=-1.960515789320$,
against the literature value $-1.960515789825$ credited to Alves et al.\ \cite{alves} (relative difference
$2.6\times10^{-10}$). The square lattice gives $e_{\mathrm{sq}}=-1.950132460$, which agrees with the corresponding literature value for
the square lattice to relative $4.3\times10^{-10}$, so $e_{\mathrm{tri}}<e_{\mathrm{sq}}$, in agreement with the classical lattice comparison of
Bonsall and Maradudin \cite{bm}. The square lattice is a saddle (64 unstable modes in the tested cell), and the
result is invariant under the Ewald splitting parameter to better than $10^{-10}$ over a factor-$4$ range. The
crate's energy convention (pairs counted once or twice, as in~\eqref{ch5:eq:jellium}) has not been matched to upstream's
here; the \emph{ordering} triangular below square does not depend on it.

Status of the main relaxation study (no relaxed periodic configuration below $e_{\mathrm{tri}}$, hypothesis E1):
\JelliumMainStatus.

Even a perfect outcome is Tier X: it would be a consistency check of Corollary~\ref{ch5:cor:jellium} in one case, not
evidence for it.
\end{numbox}

\subsection{The endpoint \texorpdfstring{$s=2$}{s=2}}

The endpoint $s=d=2$ is open, and it lies outside upstream's statement. Here is precisely what breaks.
\begin{enumerate}[leftmargin=*]
\item \textbf{The background is not integrable.} By~\eqref{ch5:eq:locint}, $\int_{K_R}|a-x|^{-2}\,\mathrm{d}x=+\infty$ for
  every $a\in K_R$, and so is the background self-energy. In~\eqref{ch5:eq:jellium} the second and third terms are both
  infinite and enter with opposite signs: $J_{2,R}$ is of the form $\infty-\infty$ and has no definition. This is the
  basic obstruction; the next three are what it looks like inside upstream's proof.
\item \textbf{The field class is empty.} The admissible range $1<p<\min(2,2/s,3/(s+1))$ becomes $1<p<1$ at $s=2$.
\item \textbf{The weight leaves the good class.} The extension weight $|y|^{\gamma}$ has $\gamma=s-1=1$, outside the
  range $-1<\gamma<1$ used for the weighted minimum principle in the confinement step.
\item \textbf{The fractional Laplacian degenerates.} In the proof of Lemma~\ref{ch5:lem:heat} the operator is
  $(-\Delta)^{1-s/2}$; at $s=2$ it is the identity, which has no Green function, and $q_s$ has a pole from
  $\Gamma(1-s/2)$.
\end{enumerate}
Items 2--4 are \tierC{} observations read off the formulas of \cite[\S8]{oai}, not theorems about every possible
approach. They do not say that $s=2$ is meaningless. They say that a different renormalisation would be needed.

\begin{remark}[a possible periodic quantity at $s=2$]\label{ch5:rem:s2}
\tierC{} (unverified sketch). The bounds that make the integral in Lemma~\ref{ch5:lem:heat} converge,
$|H_C(t)-H_D(t)|\le Ct^{-1}e^{-c/t}$ as $t\downarrow0$ and $\le Ce^{-\lambda t}$ as $t\to\infty$, do not involve $s$.
The prefactor $\kappa_sq_s/\Gamma(1-s/2)=\kappa_s2^{2-s}\pi/\Gamma(s/2)$ also has a finite limit at $s=2$. So for
periodic configurations the right-hand side
$\int_0^\infty(H_C-H_\Lat)(t)\,t^{-1}\,\mathrm{d}t$ is a finite number, and it is nonnegative by the Gaussian theorem.
What is missing is a left-hand side: an independently defined $s=2$ renormalised energy whose differences this
integral computes, and an infinite-system theory that reduces to it. We have not checked whether such a quantity
appears in the literature. This is Project P2 in \S\ref{ch5:sec:projects}, and it should not be quoted as a result.
\end{remark}

%=====================================================================
\section{The strong (per-potential) Cohn--Kumar conjecture}
\label{ch5:sec:ck}

\subsection{Statement}

Cohn and Kumar conjectured that the triangular, $E_8$ and Leech lattices are universally optimal
\cite[Conjecture~9.4]{ck}. The form of the conjecture is stronger than an energy inequality: it asks for a
\emph{certificate} for each potential.

\begin{definition}[Sharp auxiliary function]\label{ch5:def:sharp}
Let $g$ be completely monotone on $(0,\infty)$ with sufficient decay. A \emph{sharp auxiliary function} for $g$ and
$\Lat$ is a radial function $f:\R^2\to\R$ in an admissible class (integrable, with integrable Fourier transform and
enough decay of both for Poisson summation over $\Lat$ and $\Lat^*$; see \cite[Conjecture~9.4]{ck} and its context for the exact class) such
that
\begin{gather*}
f(x)\le g(|x|^2)\ \ (x\ne0),\qquad \hat f(\xi)\ge0\ \ (\xi\in\R^2),\\
f(a)=g(|a|^2)\ \ (a\in\Lat\setminus\{0\}),\qquad \hat f(w)=0\ \ (w\in\Lat^*\setminus\{0\}).
\end{gather*}
\end{definition}

The \emph{strong Cohn--Kumar conjecture} in the plane is the statement: \emph{for each} completely monotone $g$ with
sufficient decay, a sharp auxiliary function exists. \tierL{} In upstream's words, the conjecture ``asks for a sharp
auxiliary function for each sufficiently decaying completely monotone squared-distance interaction: a Fourier lower
bound that attains the lattice energy and therefore proves the comparison for periodic competitors''
\cite[\S1]{oai}.

Why the conditions certify optimality is the Cohn--Elkies-type argument of Chapter~2: the two inequalities give
$E_g(C)\ge\hat f(0)-f(0)$ for every competitor of density one, and Poisson summation on $\Lat$ with both contact
conditions gives
\[
f(0)+\sum_{a\ne0}g(|a|^2)=\sum_{a\in\Lat}f(a)=\sum_{w\in\Lat^*}\hat f(w)=\hat f(0),
\]
so the bound equals the lattice energy.

\begin{warnbox}[(status)]
The strong per-potential conjecture is \textbf{open} in the plane. Upstream proves the energy inequality for every
completely monotone $g$, and constructs sharp auxiliary functions $f_\alpha$ for every Gaussian
$G_\alpha(x)=e^{-\pi\alpha|x|^2}$. It states \tierL{}: ``Integrating the resulting energy inequalities does not
construct a sharp auxiliary for each mixed potential in the class required by Cohn and Kumar's conjecture. That
stronger per-potential quantifier is not claimed here.'' \cite[\S1]{oai}
\end{warnbox}

In dimensions $8$ and $24$, Cohn, Kumar, Miller, Radchenko and Viazovska proved universal optimality and constructed
sharp Gaussian bounds and Fourier interpolation formulas \cite[Theorems~1.4, 1.7 and 1.9]{ckmrv}. We say no more
about the per-potential form there than that.

\subsection{What integrating the Gaussians does and does not give}

It is tempting to think that the per-potential statement follows at once. Write a completely monotone $g$ (with no
constant part: by upstream's lemma \lean{energy_one}, a constant has infinite energy at density one) as
\[
g(t)=\int_{(0,\infty)}e^{-\pi\alpha t}\,\mathrm{d}\mu(\alpha),
\]
with $\mu$ a positive measure, and try
\begin{equation}\label{ch5:eq:mixture}
f(x)=\int f_\alpha(x)\,\mathrm{d}\mu(\alpha).
\end{equation}
\emph{Formally}, every condition of Definition~\ref{ch5:def:sharp} holds by linearity: $f\le g$ integrates
$f_\alpha\le G_\alpha$; $\hat f=\int\hat f_\alpha\,\mathrm{d}\mu\ge0$; at $a\in\Lat\setminus\{0\}$,
$f(a)=\int G_\alpha(a)\,\mathrm{d}\mu=g(|a|^2)$; at $w\in\Lat^*\setminus\{0\}$, $\hat f(w)=0$. So the issue is not that
mixing destroys the contact conditions or the signs. It does not.

What upstream actually integrates is different. It integrates the \emph{scalar} inequalities
$E_{G_\alpha}(C)\ge\theta_\Lat(\alpha):=\sum_{a\ne0}e^{-\pi\alpha|a|^2}$, which are explicit and measurable in
$\alpha$ on both sides; Tonelli and the lower-limit argument of upstream's \S5 then give the energy inequality for $g$.
No function is ever mixed. To turn the formal computation~\eqref{ch5:eq:mixture} into a proof, one would need at least
the following \tierC{}:
\begin{enumerate}[leftmargin=*]
\item \textbf{Measurability.} The construction is carried out for each $\alpha$ separately, with finite reference
  lists and certified inequalities chosen per range of parameters. One needs a choice $\alpha\mapsto f_\alpha$ for
  which $(\alpha,x)\mapsto f_\alpha(x)$ is jointly measurable.
\item \textbf{Absolute integrability.} To exchange Fourier transform and $\mu$-integral (Fubini), one needs
  $\int\|f_\alpha\|_{L^1}\,\mathrm{d}\mu(\alpha)<\infty$, and similarly for $\hat f_\alpha$.
\item \textbf{The admissible class.} The mixture must again have the decay that Poisson summation and the class of
  \cite[Conjecture~9.4]{ck} require. This needs bounds on the decay constants of $f_\alpha$ and $\hat f_\alpha$ that are uniform,
  or at least $\mu$-integrable, in $\alpha$.
\item \textbf{Finite values at the origin.} The quantities $f(0)=\int f_\alpha(0)\,\mathrm{d}\mu$ and
  $\hat f(0)=\int\hat f_\alpha(0)\,\mathrm{d}\mu$ must be finite separately, not only their difference.
\end{enumerate}
Upstream gives none of these estimates, and neither does this book.

\paragraph{A Riesz example shows where each end matters.}
\tierC{} For $g(t)=t^{-s/2}$, the Gamma integral gives $t^{-s/2}=\frac{\pi^{s/2}}{\Gamma(s/2)}\int_0^\infty
\alpha^{s/2-1}e^{-\pi\alpha t}\,\mathrm{d}\alpha$, so $\mathrm{d}\mu\propto\alpha^{s/2-1}\,\mathrm{d}\alpha$.
\begin{itemize}[leftmargin=*]
\item As $\alpha\to0$: Poisson summation gives $\theta_\Lat(\alpha)=\alpha^{-1}\sum_{w\in\Lat^*}e^{-\pi|w|^2/\alpha}-1
  \sim\alpha^{-1}$, so $\int_0\theta_\Lat\,\mathrm{d}\mu<\infty$ iff $s>2$. This is the finiteness of the lattice
  energy (D1b), and it is the same at the level of functions, since $\hat f_\alpha(0)-f_\alpha(0)=\theta_\Lat(\alpha)$.
\item As $\alpha\to\infty$: $\mu$ has infinite mass ($\int^\infty\alpha^{s/2-1}\mathrm{d}\alpha=\infty$ for $s>0$). The
  constraints give only $-\theta_\Lat(\alpha)\le f_\alpha(0)\le1$: the upper bound from $f_\alpha\le G_\alpha$, the lower
  from $\hat f_\alpha(0)=f_\alpha(0)+\theta_\Lat(\alpha)\ge0$. So $f(0)$ is finite only if $f_\alpha(0)$ decays fast
  enough as $\alpha\to\infty$. Upstream gives no such estimate. Note also that $g(|x|^2)\to\infty$ at $x=0$, so
  $f\le g$ does not bound $f(0)$.
\end{itemize}
For a $g$ whose measure $\mu$ is finite and supported in a compact interval $[\alpha_0,\alpha_1]\subset(0,\infty)$,
items~2--4 reduce to bounds that are locally uniform in $\alpha$. In that case the obstruction may well be technical,
but we have not checked that the construction satisfies such bounds, so we do not claim it. For Riesz potentials,
the large-$\alpha$ end is a genuine quantitative question about the construction, not a formality.

\paragraph{Why one would want the certificate anyway.}
\tierC{} The energy inequality is already proved, so a per-potential certificate would not change the minimum value.
It would give an explicit object that can be checked numerically for each $g$. It would match the form of the
conjecture as Cohn and Kumar posed it. It could also carry information that the integrated inequality loses: the
gap $g(|x|^2)-f(x)$ and the function $\hat f$ measure how far a competitor is from optimal, and that is the kind of
input that stability and uniqueness questions (\S\ref{ch5:sec:unique}) need.

\paragraph{Small research projects.}
(i) Take $g(t)=\int_1^2e^{-\pi\alpha t}\,\mathrm{d}\alpha$ and try to establish items~1--4 for upstream's construction
on $[1,2]$. (ii) Measure numerically how $f_\alpha(0)$ behaves for large $\alpha$, using upstream's verification
scripts. (iii) Build a sharp auxiliary function directly, by linear programming, for one non-Gaussian $g$
(Project M4). Each of these is a step, not a solution.

%=====================================================================
\section{Open versus closed disks: resolved (D4)}
\label{ch5:sec:disks}

Upstream measures density and energy in \emph{closed} disks $\overline B(0,R)$. A reviewer asked whether open disks
would change the theorem. They do not.

\begin{proposition}[D4]\label{ch5:prop:d4}
\tierC{} (human proof; Lean status below). For every locally finite $C\subset\R^2$: $C$ has density one in closed
disks if and only if it has density one in open disks; and for every $g$, the lower energy computed with open disks
equals the lower energy computed with closed disks. Hence upstream's theorem holds verbatim with open disks.
\end{proposition}

\begin{humanproof}
Local finiteness makes $C\cap\overline B(0,R)$ finite, so only finitely many norms $|x|$, $x\in C$, lie in $[0,R]$.
Hence there is $\delta\in(0,1)$ with no point of $C$ of norm in $[R-\delta,R)$, and then the open-disk point set at
radius $R$ equals the closed-disk point set at radius $R-\delta$. Likewise the closed-disk set at $R$ equals the
open-disk set at some $R+\delta'$ with $\delta'\in(0,1)$. So each disk-energy function takes the values of the
other at radii within distance $1$: $E^{\mathrm{open}}(R)=E^{\mathrm{closed}}(R-\delta_R)\ge
\inf_{r\ge R-1}E^{\mathrm{closed}}(r)$. Letting $R\to\infty$ gives $\liminf E^{\mathrm{open}}\ge\liminf
E^{\mathrm{closed}}$, and the symmetric argument gives the reverse. For the counts,
$N_{\mathrm{closed}}(R-1)\le N_{\mathrm{open}}(R)\le N_{\mathrm{closed}}(R)$, and $(R-1)^2/R^2\to1$, so the two densities
have the same limit. The full proof is in Chapter~2.
\end{humanproof}

Lean status: \DFourStatus.

What D4 does \emph{not} touch is uniqueness. Even with both disk conventions settled, whether $\Lat$ is the unique
minimiser among \emph{periodic} configurations of density one remains open in the plane.

%=====================================================================
\section{Uniqueness, temperature and other dimensions}
\label{ch5:sec:unique}

\subsection{Uniqueness}

Upstream is explicit \tierL{}: its theorem ``identifies the minimum value, not all minimizers'' \cite{oai}. Among
all locally finite configurations, uniqueness is in fact false. Our D3 \tierA{} proves that the configuration
$\Lat\setminus\{0\}$ (\lean{vacancy}) has density one and the same lower energy as $\Lat$ for every admissible $g$
(\lean{vacancy_energy}, \lean{vacancy_universally_optimal}), and that it is not an isometric image of $\Lat$
(\lean{vacancy_not_isometric_image}); hence \lean{universal_optimum_not_unique}. A single vacancy changes a
density-one average by nothing. The meaningful uniqueness questions are therefore of the following kinds \tierC{}:
\begin{itemize}[leftmargin=*]
\item \textbf{Periodic uniqueness.} Is $\Lat$ (up to isometry) the only periodic density-one minimiser for, say,
  every Gaussian? Open in the plane.
\item \textbf{Uniqueness up to negligible sets.} Is every minimiser equal to an isometric image of $\Lat$ outside a
  set of density zero? Even formulating this needs care with lower energies.
\item \textbf{Local versions.} Several restricted results are cited in \cite[\S1]{oai}: Faulhuber et al.\ (honeycomb,
  products of one-dimensional configurations, unions of two translates), Hardin--Tenpas (four-point and six-point
  periodic classes, under a rapid-decay assumption) and Lebl\'e (sufficiently small bounded displacements cannot
  lower the energy). We cite them only through \cite{oai}.
\end{itemize}

\subsection{Positive temperature}

The theorem is a statement at zero temperature: it compares energies, not free energies. Whether particles with a
completely monotone repulsion crystallise into a triangular arrangement at small positive temperature, and in what
sense of order, is not addressed by any result in this book's references. Physicists expect that two-dimensional
positional order at positive temperature is weaker than at zero temperature; we do not cite that literature, so
treat it as a pointer, not a claim \tierC{}.

The numerics of Chapter~3 show the gap between ``is the minimum'' and ``is reached'' \tierX{}. In the Yukawa study,
annealing reached the perfect lattice in only $1$ of $20$ runs per initial-condition family at $\kappa=2$, and in $0$
of $10$ and $2$ of $10$ runs at $\kappa=1$ (prediction P2 failed there); zero-temperature quenches at $\kappa=2$ ended
defected in $15$ of $20$ and $14$ of $20$ runs. Energetic optimality and kinetic accessibility are different
questions.

\subsection{Other dimensions}

\tierL{} Per CKMRV, Ventevogel and Nijboer proved universal optimality of $\Z$ in dimension $1$, and Cohn and Kumar
gave another proof; before CKMRV no universal optimum was known in any dimension above $1$, and the case of $A_2$
``in full generality is an important open problem'' \cite[\S1.2]{ckmrv}. CKMRV settled dimensions $8$ and $24$
($E_8$ and Leech) \cite{ckmrv}; upstream claims dimension $2$ \cite{oai}. Cohn and Kumar's conjecture singled out the
triangular, $E_8$ and Leech lattices \cite{ck,oai}. None of this is formalised here.

\paragraph{Dimension 3: fcc against bcc.}
\tierC{} Scale the face-centred cubic lattice $D_3=\{x\in\Z^3:x_1+x_2+x_3\text{ even}\}$ by $2^{-1/3}$ and the
body-centred cubic lattice $\Z^3\cup(\Z^3+\tfrac12(1,1,1))$ by $2^{1/3}$, so both have covolume one. With these scalings the two
lattices are exactly dual to each other, since $D_3^*=\Z^3\cup(\Z^3+\tfrac12(1,1,1))$ and $(c\Lambda)^*=c^{-1}\Lambda^*$. Poisson summation for a
covolume-one lattice $\Lambda\subset\R^3$ gives $\sum_{x\in\Lambda}e^{-\pi\alpha|x|^2}=\alpha^{-3/2}
\sum_{w\in\Lambda^*}e^{-\pi|w|^2/\alpha}$, hence
\[
\Theta_{\mathrm{fcc}}(\alpha)=\alpha^{-3/2}\,\Theta_{\mathrm{bcc}}(1/\alpha),\qquad
\Theta(\alpha):=\sum_{x\in\Lambda}e^{-\pi\alpha|x|^2},
\]
and at $\alpha=1$ the two Gaussian energies are exactly equal. Consequently, if fcc has strictly lower Gaussian energy
than bcc at some $\alpha_*>1$, then bcc has strictly lower Gaussian energy than fcc at $1/\alpha_*$ (because
$\Theta-1$ is the energy and the identity transfers the strict inequality of $\Theta$). Neither lattice can then be
universally optimal.

\begin{numbox}[(fcc against bcc at density one)]
\tierX{} Same script (\texttt{ch5\_checks.py}), summing $e^{-\pi\alpha|x|^2}$ over $x\ne0$ in the box
$\{-12,\dots,12\}^3$ before scaling:
\begin{lstlisting}[language=Python]
def theta3(kind: str, alpha: float, K: int = 12) -> float:
    rng = np.arange(-K, K + 1)
    X, Y, Zc = np.meshgrid(rng, rng, rng, indexing="ij")
    pts = np.stack([X.ravel(), Y.ravel(), Zc.ravel()], axis=1).astype(float)
    if kind == "fcc":
        pts = pts[(pts.sum(axis=1) % 2) == 0]  # D3, covolume 2
        pts = pts * 2.0 ** (-1.0 / 3.0)
    else:
        pts = np.vstack([pts, pts + 0.5])  # bcc, covolume 1/2
        pts = pts * 2.0 ** (1.0 / 3.0)
    r2 = (pts ** 2).sum(axis=1)
    r2 = r2[r2 > 0]
    return float(np.sum(np.exp(-math.pi * alpha * r2)))
\end{lstlisting}
Output:
\begin{lstlisting}
alpha= 0.50  fcc=1.8419799351  bcc=1.8408091753  fcc-bcc=+1.171e-03
alpha= 0.80  fcc=0.5182946954  bcc=0.5170433085  fcc-bcc=+1.251e-03
alpha= 1.00  fcc=0.2315362662  bcc=0.2315362662  fcc-bcc=-2.776e-17
alpha= 1.25  fcc=0.0855078280  bcc=0.0864032476  fcc-bcc=-8.954e-04
alpha= 2.00  fcc=0.0043777160  bcc=0.0047916420  fcc-bcc=-4.139e-04
\end{lstlisting}
fcc is lower for $\alpha>1$ (short range) and bcc for $\alpha<1$ (long range), with equality at $\alpha=1$ to rounding
error, as the duality predicts. The crossing is a floating-point observation and supports no claim by itself; a
student can make it rigorous with interval arithmetic (Project U2).
\end{numbox}

This shows only that neither fcc nor bcc is universally optimal. It does not show that no configuration in
$\R^3$ is universally optimal; that stronger statement remains conjectural, and in this book it is \tierC{}.

%=====================================================================
\section{Industry and engineering implications}
\label{ch5:sec:industry}

Everything in this section is \tierC{}. We give no performance numbers, because none of the references in this book
measures an engineering benefit of the theorem, and we will not invent one.

\begin{warnbox}[(what the theorem does and does not say)]
\begin{itemize}[leftmargin=*]
\item It is about \textbf{energy minimisation} for pair potentials that are completely monotone functions of
  squared distance, at fixed density, at zero temperature. It says the triangular lattice achieves the least lower
  energy.
\item It does \emph{not} say that a physical system reaches that state (Chapter~3: kinetics, defects), that it is the
  unique ground state (D3), or anything about non-pair, anisotropic, attractive or three-body interactions.
\item Packing, quantisation, covering and sampling optimisation are \emph{different functionals}. The theorem settles
  none of them directly.
\end{itemize}
\end{warnbox}

\subsection{Which interactions are covered}

A physical system is in scope if its effective pair potential is $V(r)=g(r^2)$ with $g$ completely monotone. Three
families matter.
\begin{itemize}[leftmargin=*]
\item \textbf{Screened Coulomb (Yukawa)} $V(r)=e^{-\kappa r}/r$. In $t=r^2$ this is $t^{-1/2}e^{-\kappa\sqrt t}$. Here
  $\sqrt t$ is a Bernstein function, so $e^{-\kappa\sqrt t}$ is completely monotone; $t^{-1/2}$ is completely monotone;
  and products of completely monotone functions are completely monotone. So Yukawa is covered directly. This is the
  model of Chapter~3 (shifted-force cutoffs in the simulations are a modelling choice and are not covered).
\item \textbf{Inverse powers} $r^{-s}$ with $s>2$: covered (our D1, D1b). Dipolar repulsion $r^{-3}$ is the case $s=3$.
\item \textbf{Unscreened Coulomb} $r^{-1}$ ($s=1$) and the logarithm ($s=0$): \emph{not} covered by the main theorem;
  only by the renormalised Corollary~\ref{ch5:cor:jellium} (\tierL{}, not formalised), which assumes a neutralising
  uniform background.
\end{itemize}

\subsection{Colloidal self-assembly and coatings}

Two-dimensional colloidal crystals are studied by video microscopy \cite{zahn} and in confined geometries
\cite{bubeck}. Where the effective interaction is a screened Coulomb repulsion or, for magnetic colloids in a field,
a dipolar repulsion (check the interaction model in each paper before relying on it), the theorem identifies the
triangular lattice as the minimum-energy arrangement at the given density. For monolayer coatings and colloidal
masks, this is a useful \emph{target statement}: an assembly process that ends in a non-triangular, defect-free
monolayer of purely repulsive, isotropic particles is not at its energy minimum, and something (kinetics, the
substrate, anisotropy, attraction) is responsible. The Chapter~3 numerics are a warning in the other direction:
defects and incomplete annealing are typical, so the theorem says what the target is, not that it is reached.

\subsection{Photonic and phononic crystals}

Many photonic and phononic devices use triangular (hexagonal) arrays of holes or rods. The design criterion there is
a band-structure property (a band gap, a dispersion), which depends on the geometry through a wave equation, not
through a pair energy. The theorem does not imply that the triangular lattice gives the best band gap. Its relevance
is indirect: if such arrays are fabricated by self-assembly of repulsive particles, the triangular arrangement is
the energetically favoured one, so the fabrication route and the design target agree. Any stronger statement needs
its own argument.

\subsection{Wigner crystals: electrons on helium, ion traps, two-dimensional materials}

Electrons on the surface of liquid helium form a two-dimensional crystal \cite{grimes}; dusty plasmas form
``plasma crystals'' \cite{thomas}. Bare electrons interact by $r^{-1}$, so the relevant mathematics is the $s=1$
jellium problem, covered only by Corollary~\ref{ch5:cor:jellium}; our numerical check above agrees with the classical
lattice comparison \cite{bm,alves}. Dusty plasmas are commonly modelled with screened (Yukawa) interactions, which
are in scope of the main theorem. Trapped-ion crystals and electron crystals in two-dimensional materials are
quantum, often confined by a non-uniform trap, and sometimes not purely two-dimensional. The theorem is classical
and assumes uniform density, so it gives a reference arrangement for the bulk of a large, nearly uniform layer at
best; edge structure and quantum fluctuations are outside it. For device engineers the value is a clean benchmark:
simulation codes for such crystals can be tested against the exact minimum (Chapter~3 used exactly that control).

\subsection{Hexagonal sampling, quantisation and sensor arrays}

Hexagonal grids are used in image sensors, antenna arrays and sampling of two-dimensional fields. The reasons there
are of three kinds, and it matters which.
\begin{itemize}[leftmargin=*]
\item \textbf{Packing.} As $\alpha\to\infty$, a Gaussian energy is dominated by the shortest distances. A heuristic
  argument therefore relates universal optimality to optimal packing: a density-one configuration with larger
  minimal distance than $\Lat$ would beat it for large $\alpha$. In the plane the optimal packing is classical and
  is not new information.
\item \textbf{Sampling.} The efficiency of a sampling lattice for an isotropic band-limited field is a packing-type
  statement about the \emph{dual} lattice in frequency space. Grating lobes of an antenna array are likewise governed
  by the dual lattice. The duality $\Lat\leftrightarrow\Lat^*$ appears in the theorem's proof too, but the theorem is not
  a sampling theorem.
\item \textbf{Quantisation.} The mean-squared quantisation error of a lattice is the normalised second moment of its
  Voronoi cell: a single-cell integral, not a pair energy. The theorem says nothing about it.
\end{itemize}
The honest summary: the triangular lattice is optimal for several unrelated planar criteria, and universal
optimality adds one more, the strongest one for pair energies. It does not imply the others, and an engineer
choosing a hexagonal array should justify it by the criterion that actually matters for the device.

\subsection{Verified numerical software for safety-critical simulation}

The most transferable lesson of this project is methodological. Its pipeline has four layers.
\begin{enumerate}[leftmargin=*]
\item A theorem is checked by the Lean kernel against stated definitions, with an axiom audit. The audit flags
  planted controls (Chapter~2). This is Tier A for the formal statement only. Statement adequacy, meaning whether the
  Lean definitions say what the prose means, is a separate human audit. In this project it has not been done: 6
  ledger flags are open.
\item Inequalities used by the proof are certified with ball arithmetic (upstream's verification scripts use
  python-flint).
\item Floating-point simulations (CVODE in rusty-SUNDIALS \cite{rusty}) are preregistered and carry positive and
  negative controls: the perfect lattice, the square-lattice saddle, incommensurate boxes (Chapter~3).
\item Each claim carries a tier, so that a reader can see what supports it \cite{elenchus}.
\end{enumerate}
For safety-critical simulation (structural, nuclear, aerospace, medical), the same separation is useful. Prove the
properties that can be proved. Certify the finite computations that can be certified. Label the rest as evidence
of a weaker kind, and test the simulation against exactly known solutions, of which a universally optimal lattice is
a clean example. We make no claim that this pipeline is ready for certification regimes; that would need its own
assessment.

\subsection{Fluid dynamics: vortex arrays, mixing and superfluids}

Point vortices in two-dimensional ideal flow interact through the logarithm. Experiments have observed relaxation of
two-dimensional turbulence in magnetised electron columns to ``vortex crystals'' \cite{fine}, and triangular vortex
lattices in rotating Bose--Einstein condensates \cite{aboshaeer}. The mathematics here is the $s=0$ case: in a
rotating frame, the uniform rotation plays a role analogous to the neutralising background, and the relevant
statement is the logarithmic part of Corollary~\ref{ch5:cor:jellium} (\tierL{}, not formalised). The analogy is a
modelling statement \tierC{}: real systems are finite, have boundaries, and (for the electron-column experiments) do
not form a uniform-density lattice. For superfluid devices such as rotation sensors, the theorem gives the reference
bulk arrangement of a large vortex array; for mixing applications, where vortex arrays are a transient structure,
it is a statement about the end state of a relaxation, not about the mixing rate.

%=====================================================================
\section{Student projects}
\label{ch5:sec:projects}

Each project has a level, a concrete first step and a deliverable that can fail. ``U'' is undergraduate, ``M''
master's level, ``P'' doctoral. Respect the book's rules: tier every claim, run a positive and a negative control
before reporting a number, and never cite a Lean name you have not checked exists.

\begin{enumerate}[leftmargin=*,label=\textbf{\arabic*.}]
\item \textbf{U1. An explicit lattice-point bound.} \emph{First step:} compute the diameter $\delta$ of a fundamental
  parallelogram of $\Lat$ and make Lemma~\ref{ch5:lem:count} explicit. \emph{Deliverable:} a constant $C$ with
  $|N(r)-\pi r^2|\le C(r+1)$ for all $r\ge0$, checked exactly (integer arithmetic on $j^2+jk+k^2$) for all $r\le1000$.
  It fails if one value violates it.
\item \textbf{U2. The fcc/bcc crossing, rigorously.} \emph{First step:} prove the duality $\Theta_{\mathrm{fcc}}(\alpha)=
  \alpha^{-3/2}\Theta_{\mathrm{bcc}}(1/\alpha)$ with the scalings above. \emph{Deliverable:} an interval-arithmetic
  certificate that $\Theta_{\mathrm{fcc}}(2)<\Theta_{\mathrm{bcc}}(2)$, including a rigorous bound on the truncated
  tail. It fails if the intervals overlap.
\item \textbf{U3 / M. The $s=0$ jellium (logarithmic) Ewald energy.} \emph{First step:} implement the mean-zero
  periodic Green function of $-\Delta$ on a torus by Ewald splitting, and check that the energy does not depend on
  the splitting parameter (positive control). Choose the pair-counting convention of~\eqref{ch5:eq:jellium} and say so.
  \emph{Deliverable:} the log-jellium energies per particle of the triangular and square lattices at density one,
  with an independent second method (for example a direct Fourier-side formula you derive) agreeing to stated
  precision, and the sign of $e_{\mathrm{tri}}-e_{\mathrm{sq}}$. Corollary~\ref{ch5:cor:jellium} predicts triangular
  lower. The project fails, and the failure must be reported, if the two methods disagree or the sign is reversed.
\item \textbf{M1. Formalise the Laplace-transform step (E4).} \emph{First step:} in the pinned toolchain, find
  Mathlib's Gamma-integral lemmas. In the Mathlib checkout used by this repository's \texttt{formal/} directory,
  \lean{Real.integral_rpow_mul_exp_neg_mul_Ioi} states
  $\int_{t>0}t^{a-1}e^{-rt}\,\mathrm{d}t=(1/r)^a\,\Gamma(a)$ for $a,r>0$. Check that it exists in Mathlib
  d13f23b7 before relying on it. \emph{Deliverable:} a kernel-checked proof that
  $t^{-s/2}=\Gamma(s/2)^{-1}\int_0^\infty u^{s/2-1}e^{-tu}\,\mathrm{d}u$ for $t,s>0$, and the positivity transfer: if
  $h\ge0$ is integrable against $t^{-s/2}$, the integral is $\ge0$. Gate it as in Chapter~2: axiom audit, and a
  planted \texttt{sorry} control that the audit flags.
\item \textbf{M2. Divergence in Lean.} \emph{First step:} state \lean{latticeEnergy (riesz s)} $=\top$ for
  $0<s\le2$, the complement of D1b. \emph{Deliverable:} a kernel-checked proof, for example by bounding the sum
  below by a constant times $\sum_{n\ge1}n\cdot n^{-s}$ using shells of $\Lat$. A negative control is
  $s=3$, where the proof must fail because D1b proves finiteness.
\item \textbf{M3. The jellium energy as a Lean definition (E5).} \emph{First step:} define $J_{s,R}$
  of~\eqref{ch5:eq:jellium} for $0<s<2$ using the Bochner integral. \emph{Deliverable:} a kernel-checked non-vacuity
  lemma, namely that $J_{s,R}$ is finite at distinct points. This requires the local integrability~\eqref{ch5:eq:locint}
  in Lean. It fails if the integrability lemma cannot be established.
\item \textbf{M4. A sharp auxiliary function for one non-Gaussian $g$.} \emph{First step:} take
  $g(t)=\int_1^2e^{-\pi\alpha t}\,\mathrm{d}\alpha=(e^{-\pi t}-e^{-2\pi t})/(\pi t)$, a radial basis such as Gaussians
  times Laguerre polynomials (whose Fourier transforms are explicit), and set up a linear program: impose
  $f\le g$ and $\hat f\ge0$ on fine radial grids, impose the contact conditions on the first shells of $\Lat$ and
  $\Lat^*$, and maximise $\hat f(0)-f(0)$. Positive control: run it for a Gaussian, where the optimum must reach
  $\theta_\Lat(\alpha)$. \emph{Deliverable:} the gap between the LP optimum and $\sum_{a\ne0}g(|a|^2)$ as the basis
  grows, and a report of where it fails: which constraint is active off the grid, and which shell's contact
  condition cannot be met. This is evidence, not proof (Tier X), unless the final function is certified with ball
  arithmetic.
\item \textbf{M5. Jellium relaxation at other $s$.} \emph{First step:} reproduce the $s=1$ Ewald controls of
  \S\ref{ch5:sec:div} (triangular against literature, square-lattice saddle). \emph{Deliverable:} the same relaxation
  test (no periodic configuration below triangular) at $s=1/2$ or $s=3/2$, preregistered, with the outcome reported
  whatever it is.
\item \textbf{P1. The strong conjecture for compactly supported mixtures.} \emph{First step:} read upstream's
  construction and decide whether $\alpha\mapsto f_\alpha$ can be chosen measurably on $[\alpha_0,\alpha_1]$.
  \emph{Deliverable:} either a proof of items~1--4 of \S\ref{ch5:sec:ck} for every $\mu$ supported in a compact
  subinterval of $(0,\infty)$, hence a sharp auxiliary function for each such $g$, or an explicit point where the
  construction's bounds are not uniform in $\alpha$.
\item \textbf{P2. A periodic $s=2$ energy.} \emph{First step:} make Remark~\ref{ch5:rem:s2} rigorous, that is, prove the
  convergence of $\int_0^\infty(H_C-H_\Lat)(t)\,t^{-1}\,\mathrm{d}t$ for periodic $C$. \emph{Deliverable:} an
  identification of this integral with the difference of an independently defined periodic quantity (for example a
  finite part of the divergent lattice sums of Proposition~\ref{ch5:prop:divrate}), or a proof that no such
  identification holds.
\item \textbf{P3. The heat-trace comparison in Lean (E3).} \emph{First step:} for a periodic $C$, prove in Lean that
  upstream's lower energy \lean{energy g C} for a Gaussian $g$ equals the periodic sum $\frac1N\sum_{i,j}\sum_{v\in L}'$
  that defines $H_C$. \emph{Deliverable:} the kernel-checked inequality $H_C(t)\ge H_\Lat(t)$ for all $t>0$, derived
  from the upstream theorem. The prior recorded before any attempt is about $30\%$.
\item \textbf{P4. Periodic uniqueness for Gaussians.} \emph{First step:} use the Gaussian auxiliary functions and
  their equality cases: a periodic minimiser must have all its difference vectors where $f_\alpha=G_\alpha$.
  \emph{Deliverable:} a proof that every periodic density-one minimiser for one fixed Gaussian is an isometric image
  of $\Lat$, or a description of what the contact sets of the known $f_\alpha$ allow. It fails if the equality set of
  $f_\alpha$ is larger than $\Lat\setminus\{0\}$.
\end{enumerate}

\section*{Exercises}
\addcontentsline{toc}{section}{Exercises}

\begin{exercise}
Show that the sum $\sum_{0<|a|\le R}|a|^{-s}$ over $\Z^2$ also behaves like $2\pi\log R$ for $s=2$. Explain why the
leading term does not depend on the lattice, and why the constant term can.
\end{exercise}

\begin{exercise}
Compute $\int_{K_R}|a-x|^{-s}\,\mathrm{d}x$ for $a$ at the centre of $K_R$, show that it is finite for $s<2$, and
find its behaviour as $s\uparrow2$.
\end{exercise}

\begin{exercise}
Verify that $q_s\kappa_s/\Gamma(1-s/2)=\kappa_s2^{2-s}\pi/\Gamma(s/2)$, and use the Gamma integral to show
$\int_0^\infty(4\pi t)^{-1}e^{-r^2/(4t)}t^{-s/2}\,\mathrm{d}t=(4\pi)^{-1}\Gamma(s/2)\,4^{s/2}r^{-s}$ for $r,s>0$.
\end{exercise}

\begin{exercise}
Prove that $t\mapsto t^{-1/2}e^{-\kappa\sqrt t}$ is completely monotone on $(0,\infty)$ directly from the definition,
for $\kappa>0$.
\end{exercise}

\begin{exercise}
In the proof of Proposition~\ref{ch5:prop:d4}, where exactly is local finiteness used? Describe what goes wrong in the
choice of $\delta$ if infinitely many points of $C$ accumulate at the circle $|x|=R$ from inside, and recall how
upstream's \lean{diskPoints} treats an infinite intersection.
\end{exercise}

\section*{Summary}
\addcontentsline{toc}{section}{Summary}
For $0<s\le2$ the raw energies are infinite: at $s=2$ the partial sums grow like $2\pi\log R$. The renormalised and
jellium energies for $0\le s<2$ are minimised by $\Lat$ according to upstream's \S8 \tierL{}, and none of it is in
Lean. At $s=2$ even the jellium energy is undefined, because $|x|^{-2}$ is not locally integrable in the plane. The
strong per-potential Cohn--Kumar conjecture is open. Mixing the Gaussian certificates satisfies its conditions
formally, but needs measurability, integrability and decay estimates that nobody has supplied. Open versus closed
disks is settled (D4). Uniqueness of the minimiser fails among all locally finite
configurations (D3) and is open among periodic ones. In engineering the theorem gives a clean target and a clean
benchmark for repulsive, isotropic pair interactions, and nothing more.
